%! AP Statistics — Unit 8, Lesson 8.4 — Homework
\documentclass[10pt]{article}
\usepackage{apstats-article}
\usepackage{apstats-boxes}
\usepackage{grading}
\gradingassignment{apstats-u08-l04-hw}
\gradingcourse{advanced_statistics}
\newcommand{\TallMath}[1]{$\displaystyle #1\rule[-1.4em]{0pt}{3.2em}$}

\begin{document}

\pageheader{Unit 8, Lesson 8.4}{Homework}
{\small\itshape Write each final answer in its box. Only what is inside a box is graded.}\par

\vspace{0.1in}
\textbf{Directions:} Show all work. Express expected counts to one decimal place when needed.

\begin{enumerate}

  \item \textbf{Car Ownership Survey.} A random sample of 240 high school students
        is classified by grade (10th / 11th / 12th) and whether they own a car (Yes / No).
        The marginal totals are shown below.

        \vspace{0.1in}
        \renewcommand{\arraystretch}{1.6}
        \begin{tabular}{lcccc}
                     & \textbf{Yes} & \textbf{No} & \textbf{Row Total} \\
            \hline
            10th grade & \makebox[2cm][r]{\smash{\raisebox{-2pt}{\gradedblank[id=hw1a, type=short, points=1, answer=24, check=numeric, tolerance=0]{0.8cm}{0.18in}}}} & \makebox[2cm][r]{\smash{\raisebox{-2pt}{\gradedblank[id=hw1b, type=short, points=1, answer=36, check=numeric, tolerance=0]{0.8cm}{0.18in}}}} & 60 \\
            11th grade & \makebox[2cm][r]{\smash{\raisebox{-2pt}{\gradedblank[id=hw1c, type=short, points=1, answer=36, check=numeric, tolerance=0]{0.8cm}{0.18in}}}} & \makebox[2cm][r]{\smash{\raisebox{-2pt}{\gradedblank[id=hw1d, type=short, points=1, answer=54, check=numeric, tolerance=0]{0.8cm}{0.18in}}}} & 90 \\
            12th grade & \makebox[2cm][r]{\smash{\raisebox{-2pt}{\gradedblank[id=hw1e, type=short, points=1, answer=36, check=numeric, tolerance=0]{0.8cm}{0.18in}}}} & \makebox[2cm][r]{\smash{\raisebox{-2pt}{\gradedblank[id=hw1f, type=short, points=1, answer=54, check=numeric, tolerance=0]{0.8cm}{0.18in}}}} & 90 \\
            \hline
            Col Total  & 96 & 144 & 240 \\
        \end{tabular}

        \vspace{0.1in}
        Compute all 6 expected counts. Show work for at least 2 cells.

        \writelines{4}

  \item \textbf{True or False.} ``If all row totals in a two-way table are equal, then all
        expected counts in the table must be equal.'' Write T or F in the T/F box:

        \vspace{0.1in}
        \rule{0pt}{0.45in}\textbf{TRUE} \hspace{2cm} \textbf{FALSE}\hspace{0.6in}\textit{T/F:}~\smash{\gradedblank[id=hw2a, type=choice, points=1, answer=F]{0.4in}{0.2in}}

        \vspace{0.1in}
        Explain your reasoning in 1--2 sentences.

        \gradedhide{\gradedhere{\gradedblank[id=hw2b, type=rubric, points=2, display={False: $E = (RT)(CT)/N$ also depends on the column totals; equal row totals make $E$ proportional to the column totals, not all equal.}]{\linewidth+6pt}{\gradedlinesheight{3}}}\writelines{9}}

  \item \textbf{Checking the Large-Counts Condition.} A two-way table has 3 rows and
        4 columns with a table total of 300. The large-counts condition is met if all
        expected counts exceed \hphantom{\blank{1cm}}.\hskip 0pt plus 1fill\smash{\gradedblank[id=hw3, type=short, points=1, answer=5, check=numeric, tolerance=0]{1cm}{0.2in}}

        \vspace{0.1in}
        Suppose all row totals equal 100 and all column totals equal 75.

        \begin{enumerate}[label=\alph*.]
            \item What is the expected count for every cell? Show the computation.

                  \gradedhide{\writelines{2}\gradedright{$E = $~\gradedblank[id=hw3a, type=short, points=1, answer=25, check=numeric, tolerance=0]{0.9in}{0.2in}}}

            \boxguard[16]
            \item Is the large-counts condition met? Justify.

                  \gradedhide{\writeline\\\gradedhere{\gradedblank[id=hw3b, type=rubric, points=2, display={Yes: every expected count is $25 \ge 5$.}]{\linewidth+6pt}{\gradedlinesheight{1}}}\writelines{9}}
        \end{enumerate}

\end{enumerate}

\begin{reflectionbox}
\textbf{Reflection:} In your own words, explain why expected counts represent what we
would see ``if the two variables were independent.'' Use the idea of row proportions
and column totals in your answer.

\writelines{3}
\end{reflectionbox}

\end{document}
